\chapter{Conclusions and Recommendations}
\label{chap:regression-conclusions}

\section{Main Findings}
\label{sec:reg-main-findings}

\subsection{Findings for Linear Regression}
Linear regression represents the most stable application for synthetic data generators. Parametric methods (Normal) achieve near-perfect CIO in continuous scenarios, while CART remains robust but introduces slightly larger confidence intervals. The number of predictors and correlation are the main sources of utility degradation.

\subsection{Findings for Logistic Regression}
Logistic regression is significantly more sensitive to synthesis artifacts. Small sample sizes ($N=100$) lead to low CIO and high bias, regardless of the generator. CART, in particular, tends to underestimate uncertainty (narrower CIs) in logistic outcomes, which is a high-risk failure mode for inferential research.

\subsection{Findings that are Consistent Across Both Models}
\begin{itemize}
    \item \textbf{Dimensionality:} Both models suffer as the number of predictors ($p$) increases.
    \item \textbf{Correlation:} High correlation ($\rho$) always degrades inferential fidelity.
    \item \textbf{Method Superiority:} In all parametric-compatible scenarios, parametric synthesis (Normal, LogReg) consistently outperforms non-parametric (CART) for inference.
\end{itemize}

\section{Comparison with Existing Literature}
Our findings align with and extend the work of researchers who have evaluated synthetic data in real-world settings (e.g., \cite{nowok2016, raab2016}). While those studies often highlight CART's flexibility for preserving complex univariate distributions, our systematic simulation shows that for \textbf{inferential quantities} (regression coefficients), CART's flexibility can lead to structural artifacts that degrade CIO. This confirms that predictive utility does not always translate to inferential utility.

Specifically, compared to the "baseline" study using real data from the UK longitudinal studies, our simulations provide a clearer picture of the \emph{mechanisms} of failure. We show that "synthetic bias" is not just a result of poor marginal matching but an inherent difficulty in capturing multi-dimensional conditional probabilities in sparse ($N/p$) regimes.

\section{Practical Recommendations}
\label{sec:reg-recommendations}

\subsection{Which Generators Perform Best Overall}
For researchers whose data aligns with standard parametric assumptions (linear, normal, logistic), \textbf{Parametric Generators} are strongly recommended over CART for preserving inferential results.

\subsection{Which Generators Perform Best Under Specific Conditions}
\begin{itemize}
    \item \textbf{High N, Low p:} CART and Parametric are both viable.
    \item \textbf{Small N:} Use Parametric, as CART is highly unstable.
    \item \textbf{Binary Outcomes:} Use LogReg synthesis; avoid CART unless the predictors have highly non-linear interactions.
\end{itemize}

\subsection{Guidance for Researchers Choosing a Generator}
Researchers should prioritize generators that match their intended analytical model. If a linear regression is planned for the final analysis, a linear (Normal) synthesis arm should be used.

\section{Limitations}
\label{sec:reg-limitations}
\begin{itemize}
    \item \textbf{Simulation Focus:} All datasets were generated from known distributions. Real-world non-linearities or "messy" data may change the relative performance of CART.
    \item \textbf{Metric Focus:} We focused on CIO. Other metrics (e.g., coverage, MSE) might reveal different tradeoffs.
\end{itemize}

\section{Future Work}
\label{sec:reg-future-work}
Future research should explore the intersection of \textbf{Differential Privacy} and inferential utility \cite{dwork2006, jordon2018}. Understanding how the injection of noise (e.g., $\epsilon$-privacy) further degrades CIO across different generators is a critical next step for the field. Expanding this research into settings combining differential privacy guarantees with rigorous disclosure analysis \cite{abowd2003} and comprehensive synthetic dataset evaluation \cite{bellovin2019, bowen2020, hsu2021} will provide practitioners with more robust frameworks. Finally, comparing these synthesis approaches against traditional multiple imputation techniques \cite{rubin1987, raghunathan2003} could yield further insights into the trade-offs between privacy and statistical validity.
